% Current v18 frontier. Stable IDs, exact graph counts, no existence assertion.
\subsection{残る八型と，再開する具体的な入口}
\label{ns18:support-agenda}

現在の高分母 support の必要候補は\emph{八型・二百四十六標識}である．
三高分母の完全 matching 定理を用い，旧十一型から
$G_{5,5}$ の六十標識，$G_{5,10}$ の三十標識，$G_{5,12}$ の二十標識を除いた．
これら三型は解決済みとして扱う．再開時には最初に
\nolinkurl{00_CONTINUATION_POLICY_ja.txt} を読む．
公開稿は構成・幾何，五次整数接着，局所置換の三篇であり，
本節の graph 障害と局所算術は五次整数接着論文で確認する．

$G_{n,j}$ の識別子は維持し，欠番を詰めない．
全候補が五頂点全てを使い，$ij$ は無向辺 $\{i,j\}$ を表す．
$\mathbf d$ は降順の次数列，$\beta_1=|E|-4$，$L$ は葉の数，
$N$ は $K_5$ 上の標識付き個数である．
\begin{center}
\footnotesize
\setlength{\tabcolsep}{3pt}
\begin{tabular}{crclrrr}
型&$|E|$&$\mathbf d$&$E$&$\beta_1$&$L$&$N$\\\hline
$G_{5,7}$&6&$(33321)$&$\{01,02,03,12,13,24\}$&2&1&60\\
$G_{5,9}$&6&$(43221)$&$\{01,02,03,04,12,13\}$&2&1&60\\
$G_{5,11}$&7&$(43322)$&$\{01,02,03,04,12,13,24\}$&3&0&60\\
$G_{5,13}$&7&$(44222)$&$\{01,02,03,04,12,13,14\}$&3&0&10\\
$G_{5,14}$&8&$(43333)$&$K_5-\{14,23\}$&4&0&15\\
$G_{5,15}$&8&$(44332)$&$K_5-\{24,34\}$&4&0&30\\
$G_{5,16}$&9&$(44433)$&$K_5-\{34\}$&5&0&10\\
$G_{5,17}$&10&$(44444)$&$K_5$&6&0&1
\end{tabular}
\end{center}

全 $1024$ 部分 graph と各自己同型群を
\nolinkurl{verify_support_agenda_v18.py} で再検算した．
表と全標識は \nolinkurl{support_agenda_verification_v18.json} から再生成できる．
非空 support は全五頂点で連結であり，三角形を含み，辺数六以上，$\beta_1\ge2$ である．
一つの全低 core があれば全分母一となる．
さらに，全ての二価頂点の二近傍は高辺で結ばれる．
六辺であることと葉を持つことは同値で，その場合は $G_{5,7},G_{5,9}$ だけである．
表は整数接着や atlas の存在数・同型類数を表さない．

\paragraph{第一候補：二つの葉付き型への係数伝播の拡張．}
$G_{5,7}$ の葉は $4$，その高隣接 core は $2$ で，高分母次数は三である．
一高分母の完全標準形を葉 $4$ に置き，三高分母 core $2$ の matching と接続する．
共通低 puncture の組は
\[
 (01;4),\quad(03;4),\quad(13;4),\quad(24;3)
\]
である．二価 core $3$ の高近傍 $0,1$ は高辺で結ばれるので，負の枝も保持する．

$G_{5,9}$ の葉も $4$ だが，高隣接 core $0$ は四高分母である．
三高分母 core $1$，二高分母 cores $2,3$ と組み合わせ，
共通低 punctures $(12;4),(13;4)$ を使う．
二価頂点 $2,3$ はともに高辺 $01$ を含む三角形に属するため，両者の負の枝を先に除かない．
葉から二価 core への既知の完全伝播公式は，隣接 core の次数が異なるこの二型には直接使えない．
その証明中の marked clutching，係数 $D$，全低低自由度を保持した手順を，
三高・四高の隣接 core へ拡張する．

\paragraph{第二候補：$G_{5,11},G_{5,13}$ の二高分母 cores の同時接続．}
$G_{5,11}$ の二価頂点は $3,4$ であり，それぞれ高三角形 $013,024$ に属する．
共通低 punctures $(13;4),(24;3)$ における通常 reverse の絶対値は四以下である．
二つの二高分母標準形と二つの三高分母 matching cores を，
例外接着と負の枝を含めて接続する．

$G_{5,13}$ は高辺 $01$ に三枚の三角形を付けた形で，
二価頂点 $2,3,4$ の高近傍はすべて $0,1$ である．
共通低 puncture を持つ高辺はないので，reverse 上界四をそのまま課さない．
低辺 $23,24,34$ の node 配分と，三つの二高分母 core の正負両枝を先に比較し，
最後に四高分母 cores $0,1$ の rows を復元する．

\paragraph{密な四型と局所理論．}
$G_{5,14}=K_{2,2,1}$，$G_{5,15}$，$G_{5,16}$，$G_{5,17}=K_5$ が残る．
$G_{5,15}$ には共通低 puncture $(23;4)$ があるが，二価 core $4$ の高近傍 $0,1$ は高辺で結ばれる．
ここでも負の枝を保持する．三高分母の matching 定理は該当する cores に適用し，
四高分母 core に無条件で拡張しない．
$G_{5,14}$ の dense mod 五の二模型の整数 lifting は，別の入口として残す．
高分母 support $H=\{ij:q_{ij}>1\}$ と素数 support $G_\ell=\{ij:\ell\mid q_{ij}\}$ を区別する．

三高分母の正符号 graph が完全 matching となることは解決済みである．
全高 $D$ の非負性の可否，rank 二 Schur 族の大域 compatibility，
固定低 row/column の二次復元を可変低データの族へどう適用するかは残る．
既知の固定 slice の無限族を再発見せず，標識付き二解を幾何学的に非同値だとも断定しない．

\paragraph{解決済みの型から再利用するもの．}
$K_4$，二葉三角形，$G_{5,5},G_{5,10},G_{5,12}$ は未解決候補に戻さない．
共通低 puncture の node 予算，二次式による固定低データの復元，
点と偶奇付き半直線による自由整数の完全処理を再利用する．
特に一高 hub の pair product は高 ray を含むと $Q-1$ であり，零積へ置き換えない．
局所候補が残った段階で，実面積・actual weights・integer framings・作用切片の同時条件へ進む．
研究の優先順位は数学的な面白さ，インパクト，幾何的構成への寄与に応じて柔軟に変更する．
